\documentclass[11pt,a4paper]{article}
\usepackage[utf8]{inputenc}
\usepackage[margin=1in]{geometry}
\usepackage{hyperref}
\usepackage{graphicx}
\usepackage{amsmath}
\usepackage{booktabs}

\title{RSCT Review: SUREON: A Benchmark and Vision-Language-Model for Surgical Reasoning}
\author{Swarm-It Research Discovery\\\small Automated RSCT-Based Analysis}
\date{March 09, 2026}

\begin{document}
\maketitle

\begin{abstract}
This document provides an automated review of the paper ``SUREON: A Benchmark and Vision-Language-Model for Surgical Reasoning'' from an RSCT (Representation-Solver Compatibility Testing) perspective. The paper achieved an RSCT relevance score of 33.9% and topic similarity of 42.2%.
\end{abstract}

\section{Paper Information}
\begin{itemize}
\item \textbf{Title:} SUREON: A Benchmark and Vision-Language-Model for Surgical Reasoning
\item \textbf{Authors:} Alejandra Perez, Anita Rau, Lee White, Busisiwe Mlambo, Chinedu Nwoye et al.
\item \textbf{Source:} \url{https://arxiv.org/abs/2603.06570v1}
\item \textbf{RSCT Relevance:} 33.9%
\item \textbf{Topic Match:} 42.2%
\item \textbf{Key Concepts:} multi-agent, safety
\end{itemize}

\section{Original Abstract}
Surgeons don't just see -- they interpret. When an expert observes a surgical scene, they understand not only what instrument is being used, but why it was chosen, what risk it poses, and what comes next. Current surgical AI cannot answer such questions, largely because training data that explicitly encodes surgical reasoning is immensely difficult to annotate at scale. Yet surgical video lectures already contain exactly this -- explanations of intent, rationale, and anticipation, narrated by experts for the purpose of teaching. Though inherently noisy and unstructured, these narrations encode the reasoning that surgical AI currently lacks. We introduce SUREON, a large-scale video QA dataset that systematically harvests this training signal from surgical academic videos. SUREON defines 12 question categories covering safety assessment, decision rationale, and forecasting, and uses a multi-agent pipeline to extract and structure supervision at scale. Across 134.7K clips and 170 procedure types, SUREON yields 206.8k QA pairs and an expert-validated benchmark of 354 examples. To evaluate the extent to which this supervision translates to surgical reasoning ability, we introduce two models: SureonVLM, a vision-language model adapted through supervised fine-tuning, and SureonVLM-R1, a reasoning model trained with Group Relative Policy Optimization. Both models can answer complex questions about surgery and substantially outperform larger general-domain models, exceeding 84\% accuracy on the SUREON benchmark while outperforming general-domain models on standard surgical perception tasks. Qualitative analysis of SureonVLM-R1 reveals explicit reasoning behavior, such as inferring operative intent from visual context.

\section{RSCT Analysis}
This paper presents research with 34% relevance to RSCT theory.

\textbf{Abstract Summary:} Surgeons don't just see -- they interpret. When an expert observes a surgical scene, they understand not only what instrument is being used, but why it was chosen, what risk it poses, and what comes next. Current surgical AI cannot answer such questions, largely because training data that explicitly encodes surgical reasoning is immensely difficult to annotate at scale. Yet surgical video lectures already contain exactly this -- explanations of intent, rationale, and anticipation, narrated by ex...

\textbf{RSCT Relevance:} The work touches on concepts related to multi-agent, safety, which have connections to the Representation-Solver Compatibility Testing framework.

\textbf{Further Analysis:} A detailed manual review is recommended to fully assess the implications for RSCT-based certification approaches.


\subsection*{RSCT Certification Metrics}
\begin{tabular}{ll}
\textbf{Relevance (R):} & 0.374 \\
\textbf{Spurious (S):} & 0.319 \\
\textbf{Noise (N):} & 0.307 \\
\textbf{Kappa ($\kappa$):} & 0.495 \\
\end{tabular}

The kappa score of 0.495 indicates limited representation-solver compatibility.


\section{Relevance to Swarm-It}
This paper was identified by the Swarm-It research discovery pipeline as potentially relevant to RSCT-based AI certification. The combined score of 37.2% places it in the moderate relevance category for further investigation.

\vspace{1em}
\hrule
\vspace{0.5em}
\small{Generated by Swarm-It Research Discovery | \url{https://swarms.network} | RSCT Certified}

\end{document}
